\subsection{OpenGL Display Pipeline}
\label{subsec:opengl-display}

The GL consumer thread inside \code{RenderThreadPool} dequeues completed
frames from a \code{FrameCompletionQueue} and uploads each frame as an
OpenGL texture.  Before any GL call is issued, the thread makes the
shared context current on itself.

Two rendering paths are supported.  The \emph{software path} targets
GL~1.1--compatible drivers (including the GDI Generic fallback) by tiling
the frame into power-of-two chunks that respect the
\code{GL\_MAX\_TEXTURE\_SIZE} limit---typically 1024 on GDI Generic---and
converting \code{Color16} pixels to RGBA8 before upload.  The
\emph{hardware path} uploads the entire frame as a single
non-power-of-two texture in RGBA16 format via \code{glTexImage2D},
avoiding the tiling overhead entirely.

Both paths flip the Y~axis during texture-coordinate generation to
account for the mismatch between the top-left origin used by the
frame buffer and the bottom-left origin expected by OpenGL.

% -----------------------------------------------------------------
\subsection{PNG Output and GPU Compression}
\label{subsec:png-output}

\FractalShark{} writes lossless 16-bit PNG images through the backend
selected by the resolved render algorithm. CPU algorithms retain
\code{WPngImage} encoding and the background save thread pool. GPU
algorithms use the CUDA encoder and finish encoding and writing on the
thread that initiated the save. AUTO follows the resolved algorithm's
\code{UseLocalColor} setting. Both paths share filename resolution,
extension handling, and the behavior of skipping an existing output file.

\subsubsection{Save dispatch and snapshot ownership}

\code{SaveCurrentFractal} captures the encoder backend, dimensions,
antialiasing factor, iteration type and limit, palette, and iteration
buffer in \code{PngParallelSave}. A copied snapshot owns an independent
iteration buffer. A moved snapshot installs a replacement current buffer
and returns its original storage to the spare-buffer pool after saving.
Queued CPU saves consequently remain valid across later renders or
resizes. Iteration-text saves also retain the background worker path.

Callers drain the render pool before reading the authoritative
\code{m\_CurIters} snapshot (\cref{subsec:two-render-paths}). The GPU save
path initializes renderer~0 from that snapshot and uploads its logical
iteration rows using the padded host row stride. The upload avoids any
dependency on which renderer produced the published frame. Existing
GPU coloring and antialiasing kernels then produce the final packed,
row-major \code{Color16} image on the compute stream. The device image
passes directly to a lazily allocated, reusable PNG encoder; no
uncompressed color image is downloaded for encoding.

GPU saves follow the existing GPU coloring semantics, including the
palette-rotation distinction described in \cref{subsec:coloring}.
Encoding is lossless with respect to those final colors. GPU encoding
and file-write failures return directly to the caller, whereas CPU
background failures use the existing worker diagnostics. The GPU bridge
waits for pending stream work before releasing the save snapshot and
does not retain its palette address as a host cache key.

\code{PngCompletionMode::Background} permits background completion only
for CPU saves. GPU saves are complete when the save call returns.
The corresponding CLI response and timing behavior is documented in
\cref{sec:fractalshark-cli}.

\subsubsection{Encoder interface and lifetime}

The CUDA-facing interface is declared in
\code{FractalSharkGpuLib/PngEncoder.cuh}:

\begin{lstlisting}
FractalShark::Png::GpuPngEncoder encoder;
std::vector<unsigned char> pngBytes;
const cudaError_t error = encoder.Encode(
    devicePixels, width, height, rowStrideBytes, stream, pngBytes);
\end{lstlisting}

The input is a caller-owned device allocation of native-endian 16-bit
RGBA samples with an explicit byte stride. Row padding is ignored;
the allocation must cover the requested image and remain valid until
the call returns. Alpha is straight, matching \code{WPngImage}.
Opaque images use RGB16, PNG color type~2; any non-opaque sample selects
RGBA16, color type~6. Low channel bits and transparent pixels' RGB
values are preserved. Output is non-interlaced and contains the PNG
signature followed by IHDR, IDAT, and IEND chunks
\cite{PNGThirdEdition2025}.

\code{Encode} returns \code{cudaSuccess} with completed host bytes,
or a CUDA error with an empty output vector. Host allocation exceptions
propagate. Dimension, stride, alignment, and size-overflow checks precede
kernel execution; accessible device storage remains a caller requirement.
Each encoder context binds to the device used for its first valid call.
Calls on one context must be serialized, while separate contexts and
streams may execute concurrently.

\code{GetWorkspaceBytes()} reports retained device allocation capacity,
excluding the input buffer and host output. Workspace grows as needed
and is reused for smaller images. Destruction releases the
workspace on its owning device. Encoding waits for its supplied stream
to obtain alpha information, compressed sizes, and completed bytes;
it does not call \code{cudaDevice\allowbreak Synchronize}. Workspace allocation
uses \code{cudaMalloc}/\code{cudaFree}, which can impose additional
allocation synchronization.

\subsubsection{Serialization and scanline filtering}

\Cref{fig:png-encoder-pipeline} summarizes the GPU save path.
\code{PngEncoder.cu} serializes samples in big-endian byte order.
Each scanline evaluates the PNG None, Sub, Up, Average, and Paeth
filters using six bytes per RGB16 pixel or eight per RGBA16 pixel
\cite{PNGThirdEdition2025}. The minimum-sum scoring matches the bundled
LodePNG heuristic, including its unsigned None score and
\code{255 - value} convention for negative differences. Equal scores
select the lowest filter number. Only the chosen filtered row is
materialized in the compression input.

\begin{figure}[!htbp]
\centering
\begin{tikzpicture}[
  stage/.style={draw, rounded corners, align=center, font=\small,
                text width=3.7cm, minimum height=1.1cm},
  flow/.style={-{Latex[length=2mm]}, thick},
  node distance=0.55cm and 0.6cm]
  \node[stage] (color) {Upload snapshot\\GPU coloring and AA};
  \node[stage, right=of color] (filter) {Serialize samples\\Select row filters};
  \node[stage, right=of filter] (deflate) {Regional DEFLATE\\encoding};
  \node[stage, below=of deflate] (frame) {GPU checksums\\and PNG assembly};
  \node[stage, left=of frame] (transfer) {Transfer PNG bytes\\to host};
  \node[stage, left=of transfer] (write) {Write output file\\Return save status};
  \draw[flow] (color) -- (filter);
  \draw[flow] (filter) -- (deflate);
  \draw[flow] (deflate) -- (frame);
  \draw[flow] (frame) -- (transfer);
  \draw[flow] (transfer) -- (write);
\end{tikzpicture}
\caption[GPU PNG save pipeline]{GPU PNG save pipeline. Coloring,
filtering, compression, checksums, and PNG assembly use device data.
Only small result metadata and completed PNG bytes return to the host;
the initiating thread waits for completion and writes the file.}
\label{fig:png-encoder-pipeline}
\end{figure}

\subsubsection{Regional DEFLATE and dynamic Huffman coding}

Filtered bytes are divided into independent regions of at most
32~KiB. A warp processes each region with a 4,096-entry hash table,
a greedy LZ77 parser, and warp-assisted match comparisons. Matches
reference only earlier positions within their own region, may overlap,
and have a maximum length of 258~bytes. Ordered position values and
\code{atomicMax} make hash insertion deterministic when positions
consumed by one match collide. Hash tables and intermediate buffers
reside in allocated device workspace.

Each region first emits a non-final fixed-Huffman block followed by an
empty stored block that restores byte alignment. If that representation
is at least as large as a stored block, a non-final stored block replaces
it. A second kernel constructs dynamic Huffman tables and replaces the
candidate only when the complete dynamic representation is strictly
smaller. Equal sizes retain the existing candidate. Block syntax and
length/distance coding follow DEFLATE \cite{RFC1951}.

\code{PngHuffman.cuh} implements iterative boundary package merge for
length-limited Huffman trees: 15~bits for literal/length and distance
codes, and 7~bits for the header alphabet. Stable frequency/symbol
ordering resolves ties deterministically. Empty and single-symbol
alphabets use two one-bit codes for decoder compatibility; the
end-of-block symbol always receives a frequency of one. Canonical
codes are reversed for DEFLATE bit emission.

The dynamic header trims trailing unused symbols and run-length encodes
the combined code lengths with symbols~16, 17, and~18. Candidate size
includes all header and payload bits, the end-of-block symbol, and the
following empty stored block's alignment and length fields. A lower
bound skips header construction when even a minimum header cannot
improve the candidate. Only winning regions replay the deterministic
parser; no per-byte token buffer is retained. Tree nodes, sorting
buffers, traversal stacks, tables, and header runs occupy reusable
device workspace.

For \code{inputBytes} filtered bytes in \code{regionCount} regions,
stored-block fallback bounds the complete zlib output by
\code{inputBytes + 5 * regionCount + 11} bytes. A CUB device prefix
scan computes output offsets, and a compaction kernel assembles the
regions in their original order.

\subsubsection{Checksums and PNG framing}

The image uses one zlib header, one final empty stored block, and one
Adler-32 trailer \cite{RFC1950}. IDAT boundaries do not start new
compression streams. Each IDAT payload is at most 64~KiB. The GPU
generates PNG chunk headers and CRCs, combining ordered warp segments
through scalar arithmetic over the reflected CRC polynomial. The host
receives small result metadata and the final PNG byte vector.

Compression selection is automatic and requires no public mode
argument. Stronger match searching and dictionaries spanning region
boundaries remain possible future improvements. The current independent
regions favor parallelism and bounded workspace at the cost of some
compression ratio.

\subsubsection{Measured encoding cost and output size}
\label{subsec:png-benchmarks}

\Cref{tab:png-current-benchmark,tab:png-compression-sizes} report
3840~by~2160 opaque images measured on October~4, 2026 with an NVIDIA
GeForce RTX~5090 (\code{sm\_120}), CUDA Toolkit/runtime~13.2, and
Visual Studio~2026 in Windows Release. The measurements were collected
after builds and correctness tests completed, with no concurrent build
or test workload. The benchmark procedure is given in
\cref{subsec:png-validation}.

GPU timing includes serialization, filtering, compression, checksums,
PNG assembly, stream waits, and transfer of completed bytes. Input
upload and construction of the CPU \code{WPngImage} occur outside the
measured interval. Warm GPU and CPU times are medians of three calls;
first-call timing separately includes initial workspace allocation.
Fractal rendering, palette conversion, iteration upload, and disk
writing are excluded. The CPU baseline uses
\code{PngEncodingOptions\{false, false, 32\}}. These measurements
characterize the fixtures and hardware rather than end-to-end save
latency or a timing guarantee.

\begin{table}[!htbp]
\centering
\small
\setlength{\tabcolsep}{4pt}
\caption[4K PNG encoding measurements]{4K PNG encoding measurements
after production save integration. Times are warm medians in
milliseconds; workspace is retained encoder allocation capacity in MiB.}
\label{tab:png-current-benchmark}
\begin{tabularx}{\textwidth}{Xrrrrr}
\hline
Image & GPU ms & CPU ms & GPU bytes & CPU bytes & MiB \\
\hline
Uniform      & 2.42  & 502.36  & 76,027     & 51,113     & 264.3 \\
Gradient     & 5.31  & 671.04  & 3,477,715  & 1,147,577  & 267.5 \\
Seeded noise & 42.79 & 1883.63 & 49,785,331 & 49,794,895 & 311.7 \\
Mandelbrot   & 13.12 & 598.19  & 857,686    & 692,127    & 265.0 \\
\hline
\end{tabularx}
\end{table}

The Mandelbrot fixture encodes in approximately 13.12~ms on the GPU
versus 598.19~ms on the CPU, while its GPU file is approximately
1.24~times the CPU size. The gradient fixture retains a larger size
penalty because the greedy parser and independent dictionaries limit
match discovery. Noise selects stored blocks and has nearly the same
output size as the CPU encoder. Retained workspace excludes the
63.3~MiB input image and all host allocations.

\begin{table}[!htbp]
\centering
\small
\caption[PNG compression size comparison]{PNG compression size comparison
on identical filtered 4K inputs. Dynamic coding replaces the fixed/stored
candidate only when smaller; equal sizes retain that candidate.}
\label{tab:png-compression-sizes}
\begin{tabularx}{\textwidth}{Xrrr}
\hline
Image & Fixed/stored bytes & Dynamic bytes & Reduction \\
\hline
Uniform      & 325,145    & 76,027     & 76.6\% \\
Gradient     & 5,082,472  & 3,477,715  & 31.6\% \\
Seeded noise & 49,785,331 & 49,785,331 & 0\% \\
Mandelbrot   & 1,180,922  & 857,686    & 27.4\% \\
\hline
\end{tabularx}
\end{table}

Dynamic code construction and parser replay trade encoding work for
smaller files and add approximately 45~MiB to retained 4K workspace
relative to the initial fixed/stored implementation. Timing is
diagnostic; exact decoded pixels and compression-size bounds are the
benchmark's assertions.
